\section{Stakeholders del proyecto}
\label{sec:stakeholders}

Esta sección consolida la lista de \textit{stakeholders} a partir del
prompt 01 (Gemini), enriquecida por el equipo con actores que el modelo
no incluyó (\textit{peer} entrevistador como rol activo,
reclutador/coach como consultor, administrador del sistema como rol
diferenciado del equipo de desarrollo, y comunidad académica de IHC).
La numeración resultante \sk{1}--\sk{17} es la \textbf{autoritativa} en
todo el documento y se referencia en las secciones de requerimientos
(\ref{sec:requerimientos}) y casos de uso (\ref{sec:casos}).

Clasificación: \textbf{P} = primario (interactúa directamente con la
aplicación), \textbf{S} = secundario (interactúa de forma indirecta),
\textbf{T} = terciario (afectado pero no interactúa).

\subsection{Listado completo}

\begin{longtable}{p{1.0cm} p{0.6cm} p{4.4cm} p{8.0cm}}
\toprule
\textbf{ID} & \textbf{Tipo} & \textbf{Stakeholder} & \textbf{Interés / rol respecto al sistema} \\
\midrule
\endhead

\sk{1}  & P & Estudiante de ingeniería en búsqueda de prácticas
preprofesionales (UNI) &
Interactúa con el sistema para entrenar habilidades blandas y técnicas
específicas del perfil de egreso. Su interés es obtener feedback
accionable para reducir la ansiedad y mejorar su desempeño en procesos
reales. \\
\addlinespace

\sk{2}  & P & Recién egresado en búsqueda de su primer empleo a tiempo
completo &
Utiliza la plataforma para simular entrevistas de mayor complejidad
técnica y estratégica según el mercado laboral actual. Busca validar
su capacidad de comunicación efectiva y el manejo de lenguaje corporal
frente a un evaluador basado en IA. \\
\addlinespace

\sk{3}  & P & Usuario en rol de \textit{peer} entrevistador &
El mismo estudiante (\sk{1} o \sk{2}) en rol invertido durante
\textit{mock interviews} colaborativas: formula preguntas, ancla
comentarios, califica. \\
\addlinespace

\sk{4}  & P & Usuario con discapacidad visual (baja visión o ceguera) &
Accede a la interfaz mediante lectores de pantalla y comandos de voz,
requiriendo cumplimiento de WCAG 2.2 nivel AA. Su interés es el
entrenamiento de la voz y el ritmo sin barreras de accesibilidad
visual. \\

\sk{5}  & P & Usuario con discapacidad auditiva (hipoacusia o sordera) &
Interactúa mediante subtitulado en tiempo real y señales visuales que
compensan la retroalimentación sonora. Depende de una interfaz que
traduzca el tono y ritmo del entrevistador en indicadores claros. \\

\sk{6}  & P & Usuario con discapacidad motora (limitación funcional en
extremidades superiores) &
Utiliza la aplicación mediante tecnologías de asistencia como
pulsadores o control por voz. Requiere que la interfaz no dependa de
interacciones físicas complejas o tiempos de respuesta táctiles
exigentes. \\

\sk{7}  & P & Usuario con condición de neurodivergencia (TEA, TDAH,
dislexia) &
Emplea el simulador para practicar el contacto visual y la
interpretación de señales sociales en un entorno controlado y
predecible. Su interés principal es la reducción de la sobrecarga
cognitiva mediante una interfaz clara y la repetición ilimitada de
sesiones. \\
\addlinespace

\sk{8}  & S & Reclutador o profesional de RRHH (consultor del proyecto) &
Aporta criterios reales de evaluación de entrevistas; revisa la
calibración de las preguntas y rúbricas generadas por el LLM. \\
\addlinespace

\sk{9}  & S & Coach de carrera o mentor &
Recomienda la herramienta, revisa el plan de mejora del candidato y
puede dejar comentarios asincrónicos en sesiones grabadas. \\
\addlinespace

\sk{10} & S & Oficina de Empleabilidad de la Universidad Nacional de
Ingeniería &
Monitorea reportes agregados y anonimizados para identificar brechas
de habilidades comunes en los estudiantes. Utiliza estos datos para
ajustar talleres de inserción laboral y mejorar los indicadores de
empleabilidad de la institución. \\
\addlinespace

\sk{11} & S & Equipo de desarrollo y mantenimiento del proyecto &
Responsable de la implementación técnica, mantenimiento de los modelos
de MediaPipe y la integración de las API multimodales. Asegura
estabilidad operativa del backend y persistencia de los planes de
mejora personalizados. \\
\addlinespace

\sk{12} & S & Administrador del sistema &
Gestiona usuarios, ligas, contenido de industrias y moderación del
\textit{feed} comunitario. Diferenciado del equipo de desarrollo
(\sk{11}) por tener acceso a operaciones en producción y no a
despliegues. \\
\addlinespace

\sk{13} & S & Docentes del curso CC451 (cátedra) &
Supervisan el uso del sistema como herramienta pedagógica para
evaluación de métricas de UX y accesibilidad. Validan que el proyecto
cumpla con los objetivos académicos y técnicos planteados en el ciclo
2026-I. \\
\addlinespace

\sk{14} & S & Administradores de infraestructura en la nube
(Render, Fly.io, Vercel) &
Proveen servicios de cómputo y red para el despliegue de la PWA y el
backend. Su interés es el consumo eficiente de recursos y el
cumplimiento de los acuerdos de nivel de servicio. \\
\addlinespace

\sk{15} & T & Empresas reclutadoras del sector tecnológico y de
ingeniería &
Beneficiarias indirectas: reciben candidatos mejor preparados. No
interactúan directamente con el software, pero el resultado del
entrenamiento impacta en la calidad de sus procesos de selección. \\
\addlinespace

\sk{16} & T & Autoridad Nacional de Protección de Datos Personales
(Perú, ANPDP) &
Vela por el cumplimiento de la Ley N\textsuperscript{\textordmasculine}~29733
respecto al tratamiento de datos biométricos y grabaciones procesadas
por el sistema. Su rol es regulatorio: asegurar que no se vulnere la
privacidad de los estudiantes durante el análisis de voz y video. \\
\addlinespace

\sk{17} & T & Familiares y entorno cercano del candidato &
Se ven influenciados por la mejora en las perspectivas económicas y el
éxito profesional del usuario tras el uso del simulador. Representan
el entorno de apoyo que se beneficia indirectamente de la inserción
laboral efectiva del candidato. \\

\bottomrule
\end{longtable}

\subsection{Mapa de prioridades}
Para el alcance del MVP se priorizan los \textit{stakeholders}
primarios (\sk{1}--\sk{7}), garantizando que las rutas de uso de los
cuatro perfiles con discapacidad (\sk{4}--\sk{7}) sean equivalentes en
funcionalidad a la del candidato general (\sk{1}). Los
\textit{stakeholders} secundarios y terciarios influyen en
requerimientos no funcionales (privacidad, mantenibilidad, costo de
operación, cumplimiento normativo) y se atienden en la
sección~\ref{sec:requerimientos}.

\subsection{Equivalencia con la numeración propuesta por Gemini}
La respuesta cruda del prompt 01 (versionada en
\texttt{prompts/respuestas/gemini\_01\_stakeholders.md}) usa una
numeración propia de 14 \textit{stakeholders}. La equivalencia con la
numeración autoritativa del informe es:

\begin{itemize}
  \item Gemini-S1 $=$ \sk{1}; Gemini-S2 $=$ \sk{2}.
  \item Gemini-S3, S4, S5, S6 $=$ \sk{4}, \sk{5}, \sk{6}, \sk{7}
        (perfiles con discapacidad).
  \item Gemini-S7 $=$ \sk{10} (Oficina de Empleabilidad UNI).
  \item Gemini-S8 $=$ \sk{13} (Docentes del curso).
  \item Gemini-S9 $=$ \sk{11} (equipo de desarrollo).
  \item Gemini-S10 $=$ \sk{14} (infraestructura cloud).
  \item Gemini-S11 $=$ \sk{15} (empresas reclutadoras).
  \item Gemini-S12 $=$ \sk{16} (autoridad de protección de datos).
  \item Gemini-S13 $=$ \textit{proveedor LLM externo}, no incluido como
        \textit{stakeholder} en el informe; aparece como dependencia
        técnica en arquitectura.
  \item Gemini-S14 $=$ \sk{17} (familiares).
\end{itemize}

\noindent El equipo agregó cuatro \textit{stakeholders} no propuestos
por Gemini: \sk{3} (peer entrevistador como rol activo), \sk{8}
(reclutador/RRHH como consultor del proyecto), \sk{9} (coach/mentor
con feedback asincrónico) y \sk{12} (administrador del sistema
diferenciado del equipo de desarrollo).
